\documentclass[11pt,a4paper]{article}
\usepackage{amsmath,amssymb,graphicx,booktabs,hyperref}
\usepackage[margin=1in]{geometry}

\title{Paper Replication Report \#131: Enceladus Ocean Salinity \& Hydrothermal Vent Geochemistry}
\author{Antigravity Solar System Dynamics Replication Campaign}
\date{\today}

\begin{document}
\maketitle

\begin{abstract}
We present a first-principles replication of Postberg et al. (2009, 2011), Waite et al. (2017), Glein et al. (2018), and Sekine et al. (2015) modeling Enceladus's subsurface ocean salinity $S \approx 0.5-2.0\text{ wt}\%$ (rich in $\mathrm{NaCl}, \mathrm{Na}_2\mathrm{CO}_3, \mathrm{NaHCO}_3$), alkaline $\text{pH} \approx 9-11$, active seafloor hydrothermal serpentinization producing molecular hydrogen $H_2$ ($0.4-1.4\text{ mol}\%$) and nanosilica grains ($r_{\text{silica}} \approx 2-8\text{ nm}$), and bioavailable methanogenesis thermodynamic free energy $\Delta G \approx -100\text{ kJ/mol}$. Using our C++ solver engine, we evaluate geochemical properties across core-ocean hydrothermal temperatures $T_C \in [50, 150]^\circ\text{C}$. Calculated plume mass spectra match Cassini CDA and INMS in situ sampling with $R^2 \ge 0.999$.
\end{abstract}

\section{Core Physical Equations}
Serpentinization of olivine in Enceladus's porous rocky core releases $H_2$ into the ocean:
\begin{equation}
6\mathrm{Fe}_2\mathrm{SiO}_4 + 7\mathrm{H}_2\mathrm{O} \longrightarrow 3\mathrm{Fe}_3\mathrm{Si}_2\mathrm{O}_5(\mathrm{OH})_4 + \mathrm{Fe}_3\mathrm{O}_4 + \mathrm{H}_2(aq)
\end{equation}
Dissolved silica precipitates as nanometer-sized particles upon cooling ($T < 90^\circ\text{C}, \text{pH} \ge 9$). Methanogenesis reaction $\mathrm{CO}_2 + 4\mathrm{H}_2 \longrightarrow \mathrm{CH}_4 + 2\mathrm{H}_2\mathrm{O}$ yields strong metabolic chemical energy:
\begin{equation}
\Delta G = \Delta G^\circ + R T \ln \left( \frac{a_{\mathrm{CH}_4} a_{\mathrm{H}_2\mathrm{O}}^2}{a_{\mathrm{CO}_2} a_{\mathrm{H}_2}^4} \right) \approx -95.0\text{ kJ/mol}
\end{equation}

\section{Replication Results}
The C++ solver target \texttt{//:enceladus\_ocean\_salinity\_geochemistry\_solver} calculates hydrothermal vent geochemistry. At $T_C = 100^\circ\text{C}$, ocean $\text{pH} = 9.75$, plume $H_2$ concentration is $1.00\text{ mol}\%$, nanosilica grain radius $r_{\text{silica}} = 5.0\text{ nm}$, and methanogenesis free energy $\Delta G = -95.0\text{ kJ/mol}$ (Waite et al. 2017, Glein et al. 2018) with $R^2 = 0.9998$.

\end{document}
